次に、これらの特性が高被引用論文に特有のものであるかを確認するため、表\ref{tab:g2obsindex}にG2論文の指標値を示す。
本指標値の算出には十分な被引用数が必要であるため、比較するグループはG2論文のみとする。
これらの表によれば高被引用論文G1はG2に比べ全体的に年齢が高い。メンバー構成においては、高被引用論文はG2に比べクラスGとクラスSBの構成比率が高くなっており明らかに異なる。
また、十分に成熟しオブソレッセンスの傾向がみられる群、すなわちクラスNoのピーク年（出版から被引用ピークまでの期間、中央値）を、G1とG2とで確認したところ、G1は9.5年、G2は5年であり、既存の報告（前述Paroloら）と比べるとG1では長くなる傾向を見せ、G2では整合的であった。
